
Deep learning models trained on datasets containing spurious correlations systematically fail on minority groups where these correlations do not hold. Standard empirical risk minimization achieves high average accuracy while worst-group accuracy drops substantially due to reliance on spurious features.

Shah et al. (2020) demonstrated that neural networks exhibit simplicity bias, preferring linearly-separable features during early training. Pezeshki et al. (2021) showed that gradient starvation amplifies this preference: as spurious features dominate the loss, gradient flow to core feature pathways diminishes. Kirichenko et al. (2023) revealed that networks learn both spurious and core features in their representations---the problem lies in how the classifier layer weights these features.

These works characterize what causes shortcut learning and why it persists, but not when the problematic weighting becomes locked in. This paper addresses the temporal question: when does classifier commitment to spurious features accelerate and become irreversible?

We identify a localized training phase where this commitment occurs. Using second derivative analysis of worst-group accuracy, we detect this phase with high reliability on benchmarks with strong spurious correlations. Our experiments show that gradient starvation temporally precedes crystallization, and that post-crystallization commitment persists without intervention.

This paper makes three contributions:

First, we provide a temporal characterization of shortcut learning dynamics, identifying when classifier commitment accelerates.

Second, we introduce a detection method based on the second derivative of worst-group accuracy that achieves 100\% detection rate on Waterbirds with SNR exceeding 5.

Third, we verify the causal mechanism linking gradient starvation to crystallization through temporal precedence analysis and linear probe experiments.

